การตรวจชิ้นเนื้อทางศัลยพยาธิวิทยาด้วยตาเปล่า (Gross Examination) ต้องบันทึกข้อมูลขณะสวมถุงมือที่เปื้อนสิ่งส่งตรวจและฟอร์มาลิน อีกทั้งมีเสียงพัดลมตู้ดูดควัน การพูดเร็ว และการแก้ไขข้อความกลางประโยค ทำให้ระบบแปลงเสียงเป็นข้อความทั่วไปถอดความได้ไม่แม่นยำ โครงงานนี้พัฒนาระบบ PathoWhisper ที่ทำงานแบบออฟไลน์สำหรับการตรวจชิ้นเนื้อมะเร็งเต้านม ประกอบด้วยการกรองเสียงรบกวน (afftdn) การถอดความด้วย Faster-Whisper Small แบบ INT8 ร่วมกับคลังศัพท์ CAP และตัวสกัดข้อมูลแบบกำหนดแน่นอนเพื่อกรอกฟิลด์ 15 ฟิลด์และสร้างรายงาน PDF

ผลทดสอบบนชุดข้อมูลเสียงสังเคราะห์ 1,000 กรณี 10 หมวด (หมวดเสียงตู้ดูดควัน 100 กรณีซ้อนทับเสียงพัดลมที่บันทึกจริง) เมื่อเทียบกับ Whisper Small พื้นฐาน พบว่า WER เฉลี่ยรายกรณีลดจาก 36.27\% เป็น 29.05\% (ช่วงความเชื่อมั่น 95\% ของการลด 6.28--8.16 จุด) และเมื่อปรับรูปแบบการเขียนมิติให้ตรงกันลดจาก 7.25\% เป็น 2.78\% ในหมวดเสียงตู้ดูดควัน ($-10$ dB SNR) WER ลดจาก 44.88\% เป็น 31.94\% เวลาประมวลผลเฉลี่ยลดจาก 20.04 เป็น 11.50 วินาทีต่อกรณี (เร็วขึ้น 1.74 เท่า) การสกัดข้อมูลประเมินครบ 15 ฟิลด์ โดย \nSpoken{} ฟิลด์มีการกล่าวถึงในชุดทดสอบ ค่า F1 เฉลี่ย (Macro) ของฟิลด์เหล่านี้เปลี่ยนจาก \MacroBase\% เป็น \MacroPW\% และอีก \nNeg{} ฟิลด์ไม่ถูกกล่าวถึงจึงใช้เป็นฟิลด์ควบคุมเชิงลบ ความต่างของ F1 ระหว่างสองระบบน้อยกว่าความต่างของ WER มาก ฟิลด์ที่ได้ผลต่ำที่สุดคือขนาดก้อน 3 มิติ (Recall \RPWMassDims\%) การทดสอบหน่วยของกฎแก้ไขกลางประโยคและคำปฏิเสธผ่าน \unitPass{} จาก \unitTotal{} กรณี ระบบเป็นเครื่องมือช่วยบันทึก ผู้ใช้ต้องตรวจทานก่อนยืนยันผล